\section{WASP-121\,b}
\label{sec:wasp121}

WASP-121\,b is an ultra-hot Jupiter orbiting very close to its Roche limit,
which makes its upper atmosphere a benchmark for atmospheric escape studies.
Its transmission spectrum is exceptionally metal-rich, and several atomic
features (H$\alpha$, Ca\,II, Mg\,II, Fe\,II, Ba$^{+}$) are observed to extend
beyond the transit-equivalent Roche-lobe radius ($R_{\rm eqRL}\simeq
1.3$--$1.6\,R_{\rm p}$), directly tracing the escaping outflow. All entries
below are drawn from real ADS records; bibcodes are given in
\texttt{typewriter} font.

\subsection{Observations}

The escape-tracing detections are dominated by VLT high-resolution
spectroscopy (ESPRESSO, UVES, HARPS) at optical wavelengths and HST/STIS in
the near-UV; the He\,I triplet was first mapped by JWST/NIRISS. No dedicated
Ly$\alpha$ (1215.67\,\AA) detection for WASP-121\,b appears in ADS, and there
is no published ground-based He\,I\,10830 line-detection paper prior to the
JWST phase curve: both remain observational gaps (the F-type host star
weakens the He signal).

\begingroup
\setlength{\tabcolsep}{4pt}
\renewcommand{\arraystretch}{1.15}
\begin{longtable}{@{}>{\raggedright\arraybackslash}p{2.4cm}>{\raggedright\arraybackslash}p{2.3cm}>{\raggedright\arraybackslash}p{5.0cm}>{\raggedright\arraybackslash}p{3.2cm}@{}}
\toprule
\textbf{Species / line} & \textbf{Instrument} & \textbf{Measured result} & \textbf{Reference} \\
\midrule
\endfirsthead
\toprule
\textbf{Species / line} & \textbf{Instrument} & \textbf{Measured result} & \textbf{Reference} \\
\midrule
\endhead
\bottomrule
\endfoot
Mg\,II 2796/2804, Fe\,II UV1/UV2 & HST/STIS NUV (PanCET) &
Mg\,II resolved at $5.9\sigma$, reaching $R_{\rm pl}/R_{\star}=0.284\pm0.037$
(2796\,\AA); Fe\,II multiplets to $R_{\rm pl}/R_{\star}\approx0.3$; both exceed
the transit Roche size ($R_{\rm eqRL}/R_{\star}=0.158$), i.e.\ unbound /
escaping refractories &
Sing et al.\ 2019 \texttt{2019AJ....158...91S} \\
H$\alpha$ 6562.8, Na\,I, Fe\,I & VLT (archival high-res optical) &
New H$\alpha$ detection, transit depth $1.87\pm0.11\%$, attributed to an
extended (escaping) H atmosphere; Na\,I doublet confirmed; Fe\,I via
cross-correlation (photospheric) &
Cabot et al.\ 2020 \texttt{2020MNRAS.494..363C} \\
Fe\,(atomic, limb) & HARPS (HEARTS III) &
Reloaded Rossiter--McLaughlin isolates planetary-limb metal (likely Fe\,I)
absorption; line width $14.3\pm1.2$\,km\,s$^{-1}$; blueshift $-5.2\pm0.5$\,km\,s$^{-1}$
(day--night winds or thermosphere expansion) &
Bourrier et al.\ 2020 \texttt{2020A\&A...635A.205B} \\
Fe\,I (blue arm) & VLT/UVES &
Neutral Fe detected $>8\sigma$; upper-atmosphere $T=3710^{+490}_{-510}$\,K;
Fe\,I found physically separated from exospheric Fe\,II &
Gibson et al.\ 2020 \texttt{2020MNRAS.493.2215G} \\
H$\alpha$, Ca\,II, Na, K, Li, Mg, Fe\,I/II, Cr\,I, V\,I & VLT/\allowbreak ESPRESSO &
H$\alpha$ ($\sim1.44\,R_{\rm p}$) and Ca\,II ($\sim2\,R_{\rm p}$) extend beyond
the Roche lobe with broadening inconsistent with tidally-locked rotation
alone (winds / high-altitude jets); Li detected at $6\sigma$ &
Borsa et al.\ 2021 \texttt{2021A\&A...645A..24B} \\
H\,I, Ca\,II (exospheric) + 17 atomic species & VLT/UVES &
Confirms exospheric H\,I and Ca\,II; 5 strong detections (incl.\ Cr\,I, V\,I,
Ca\,I, K\,I), Fe\,I at $8.8\sigma$, novel Sc\,II at $4.2\sigma$ &
Merritt et al.\ 2021 \texttt{2021MNRAS.506.3853M} \\
Fe, Cr, V (relative abundances) & VLT/UVES &
Retrieved abundances near-solar; anomalously large Mg attributed to the
escaping atmosphere not captured by the (bound) model &
Gibson et al.\ 2022 \texttt{2022MNRAS.512.4618G} \\
Ba$^{+}$, Sr$^{+}$, Co, Ca$^{+}$ (+ Ti$^{+}$ tent.) & VLT/\allowbreak ESPRESSO &
First Ba$^{+}$ detection (heaviest species then known); Ca$^{+}$ shows a clear
blueshifted asymmetric feature reaching beyond the Roche lobe, consistent with
atmospheric escape &
Azevedo Silva et al.\ 2022 \texttt{2022A\&A...666L..10A} \\
H$\alpha$, Fe\,II, Ca\,II (exospheric) & VLT/\allowbreak ESPRESSO (multi-epoch) &
Retrieved extended features: H$\alpha$ $1.54\pm0.04\,R_{\rm p}$, Fe\,II
$1.17\pm0.01\,R_{\rm p}$, Ca\,II $2.52\pm0.34\,R_{\rm p}$ (beyond $R_{\rm eqRL}\sim1.3\,R_{\rm p}$);
consistent across epochs/instruments &
Maguire et al.\ 2023 \texttt{2023MNRAS.519.1030M} \\
Fe, Ni, C, O (refractory vs volatile) & VLT/CRIRES$^{+}$ + ESPRESSO &
Bounded abundances of refractory Fe, Ni and volatile C, O; atmosphere enriched
in volatiles relative to refractories (photospheric composition, formation
tracer) &
Pelletier et al.\ 2025 \texttt{2025AJ....169...10P} \\
Fe, CO, V, H$_2$O & HARPS + NIRPS + CRIRES$^{+}$ &
Cross-correlation Fe/CO/V ($5.8/5.0/4.7\sigma$); $\Delta K_p=-15\pm3$\,km\,s$^{-1}$
circulation offset (photospheric composition/dynamics) &
Vaulato et al.\ 2025 \texttt{2025A\&A...703A.251V} \\
He\,I 10830 (metastable) & JWST/NIRISS (phase curve) &
Absorption persists $>17$\,h ($55\%$ of orbit), leading transit by $>5$\,h and
trailing by $>10$\,h; optically thick region $\sim110\,R_{\rm p}$ ($\sim0.1$\,AU),
resolved into leading and trailing tails; not self-consistently reproduced by
current models &
Allart et al.\ 2026 \texttt{2026ASTCS..1130205A} \\
\end{longtable}
\endgroup

\subsection{Model studies}

\paragraph{Huang, Koskinen, Lavvas et al.\ 2023 (\texttt{2023ApJ...951..123H}).}
\textit{Code / method:} 1D forward model of the upper atmosphere and
hydrodynamic escape, extended to include atomic metal species (Mg, Fe, Ca),
coupled to a Ly$\alpha$ Monte-Carlo radiative-transfer calculation for the
excited (n\,=\,2) hydrogen population and its heating/ionization feedback.
\textit{Initial conditions / setup:} 1D radial escape model driven by stellar
XUV; explicitly treats Roche-lobe overflow and requires careful coupling to
the structure of the lower and middle atmosphere (base boundary); metal species
carried in the ionization/thermal balance. \textit{Key results:} Reproduces the
observed Mg, Fe, Ca and H$\alpha$ features extending outside the Roche lobe;
line broadening from the Roche-lobe-overflow-driven outflow is required to
match observed line widths; the high mass-loss rate is set by Roche-lobe
overflow; metals and excited-state H significantly affect the thermal and
ionization balance of the thermosphere. (The most directly relevant benchmark
for EXHALE-type modeling.)

\paragraph{Yan, Guo, Huang \& Lavvas 2021 (\texttt{2021ApJ...907L..47Y}): BOTH (model + data comparison).}
\textit{Code / method:} XUV-driven 1D hydrodynamic escape simulation coupled to
detailed hydrogen level-population and radiative-transfer calculations for
H$\alpha$. \textit{Initial conditions / setup:} Fiducial XUV irradiation with a
scaling of $0.5$--$1.5\times$ to probe stellar-activity variation; 1D transonic
planetary wind. \textit{Key results:} Mass-loss rate $\dot{M}\sim1.28\times10^{12}$\,g\,s$^{-1}$;
high temperature and Ly$\alpha$ intensity populate the H first excited state so
the modeled H$\alpha$ transmission matches observations; the supersonic
(transonic) wind region contributes most of the H$\alpha$ equivalent width,
making H$\alpha$ a good tracer of the escaping atmosphere.

\paragraph{Wang, Lin, Wang et al.\ 2026 (\texttt{2026arXiv260201364W}).}
\textit{Code / method:} 3D simulations from the planetary surface to extended
outflows using the GPU-accelerated \textsc{Kratos} framework, coupling
hydrodynamics with non-equilibrium thermochemistry and ray-tracing radiative
transfer. \textit{Initial conditions / setup:} Full 3D domain including
day--night circulation of the lower atmosphere; transonic photoevaporative
outflow at altitude; parametric variation of FUV, EUV and X-ray irradiation;
optional stellar-wind confinement. \textit{Key results:} Outflow is shaped by
stellar gravity and orbital motion into two spiral arms with line-of-sight
velocities $\sim40$\,km\,s$^{-1}$, explaining high-velocity Na and H$\alpha$
absorption without a super-rotation jet; different species trace different
regions, Fe (inner, rotation-dominated), Na (dense spiral arms), H$\alpha$
and He\,10830 (extended, ionized gas); stellar-wind confinement enhances
metastable-He absorption.

\paragraph{Lothringer, Fu, Sing et al.\ 2020 (\texttt{2020ApJ...898L..14L}).}
\textit{Code / method:} PHOENIX 1D atmosphere models (chemical-equilibrium,
\emph{not} an escape model). \textit{Initial conditions / setup:} Grid of
equilibrium temperatures $1000$--$4000$\,K, with and without condensate
rainout. \textit{Key results:} Near-UV/blue absorption shortward of $0.5\,\mu$m
in ultra-hot Jupiters (incl.\ WASP-121\,b) can arise from equilibrium atomic/ionic
opacity (Fe, Ti, Ni, Ca, Cr, Mn, SiO) rather than only escaping metals; provides
a competing (bound-atmosphere) interpretation of the NUV excess and a rainout
diagnostic.

\paragraph{Gebek \& Oza 2020 (\texttt{2020MNRAS.497.5271G}).}
\textit{Code / method:} Non-hydrostatic exosphere / evaporative
transmission-spectrum modeling (alkaline Na\,I, K\,I curve-of-growth).
\textit{Initial conditions / setup:} Three non-hydrostatic density profiles:
escaping, exomoon, and torus. \textit{Key results:} Na\,I equivalent widths of
$\sim1$--$10$\,m\AA{} are driven by evaporation rates $\sim10^{3}$--$10^{5}$\,kg\,s$^{-1}$;
a toroidal (non-hydrostatic) Na configuration is suggested to be consistent with
the WASP-121\,b Na\,I data, breaking the atmosphere--exosphere degeneracy via the
D2/D1 line ratio.
